\section{Gaussian insertion by potential permutations}
\label{sec:insertion}

Fix a microscopic subsequential law \(\nu=\sum_j\delta_{x_j}\)
from Sections~\ref{sec:entropy}--\ref{sec:resolvents}. We shall approximate
its law by eigenvalues of a finite diagonal matrix plus Gaussian noise.
The number of potential changes will grow only inside this already
constructed limiting probability space. Thus the finite-list Gaussian
mark theorem is sufficient; no assertion about a growing list of finite-graph
vertices is needed.

The mechanism is to make many small, measure-preserving potential
changes. The potential-switch identity will represent their effect,
after truncating the spectral sum, by signed rank-one perturbations of
a diagonal matrix of microscopic positions. Conditional on those positions
and the sampled environments, the Gaussian marks make the perturbation's
mean scalar and its covariance that of a GOE matrix. We will compare it
with Gaussian noise after controlling its spectral tails. The scalar
mean need not vanish, so it will remain in the comparison matrix.

The matrix estimates below adapt the Gaussian insertion argument of
\cite[Sections~5.3--5.5]{bu}. We give their proofs because the coefficients
here come from potential permutations, and their size and cancellation
must be established for the disordered model.

\subsection{Potential permutations and their coefficients}

We use \(r\) sampled roots and potential increments of size comparable
to \(\theta\), chosen so that their total variance has scale
\(T=r\theta^2\). Two spectral cutoffs serve different purposes.
A large outer cutoff \(L\) makes the switch determinant a finite matrix
expression, with a negligible tail even though the determinant has up to
\(r\) coordinates. We will then remove the positions outside the much
smaller window \([-K,K]\). The remaining matrix has order \(K\), small
enough for the Gaussian replacement error to vanish.

Decrease the deterministic discrepancy exponent in Lemma~\ref{lem:discrepancy}
if necessary so that \(0<\chi<1/4\), and fix the following scales:
\begin{equation}\label{eq:sw-parameters}
 0<\tau<\chi/100000,\qquad
 T=r^\tau,\quad K=r^{1/100},\quad L=r^{100/\chi},\qquad
 \theta=(T/r)^{1/2}.
\end{equation}
All limits in this section have \(r\to\infty\). For clarity, the
almost-sure consequences used later will be taken along the deterministic
sequence \(r=2^j\); any faster increasing integer sequence would also do.

At each independently sampled tree root \(\alpha\), write
\[
 m_\alpha^*(v)=(v-E-\textstyle\sum_{i=1}^d g_{\alpha i})^{-1},\qquad
 a_\alpha(v)=\Re m_\alpha^*(v),\qquad
 c_\alpha(v)=\frac{\Im m_\alpha^*(v)}{\pi\rho}.
\]
Here \(v\in[-w,w]\) is the root potential and the incoming cavity
values depend only on the environment away from that root. They are finite
with positive imaginary parts almost surely. We can therefore choose a
fixed event \(\mathcal A\), measurable without the root potential, of
positive probability, on which
\[
 |E+\textstyle\sum_i\Re g_{\alpha i}|\le C_0,\qquad
 c_0\le\textstyle\sum_i\Im g_{\alpha i}\le C_0
\]
for fixed \(c_0,C_0>0\). On \(\mathcal A\), the functions \(a_\alpha,c_\alpha\)
are uniformly bounded and Lipschitz on \([-w,w]\), and
\(c_\alpha\) is bounded below by a positive constant. These statements
follow directly by differentiating \((v-A-iB)^{-1}\) on the displayed
compact range of \(A,B,v\).

Choose an integer \(N_r\) such that \(\ell_r=w/N_r\asymp\theta\), and
partition \([-w,w]\) into \(N_r\) intervals of length \(2\ell_r\).
On \(\mathcal A\), interchange the two halves of each interval by
translation. Off \(\mathcal A\), leave the potential unchanged. This is a
measure-preserving involution conditionally on the root-excluded
information. Denote its increment at the first \(r\) sampled roots by
\(D_\alpha\); at an active root, \(D_\alpha=\pm\ell_r\).
Proposition~\ref{prop:potential-switch} applies to these choices.

Evaluate \(a_\alpha,c_\alpha\) at the original root potential, and define
\begin{equation}\label{eq:sw-coefficient-bounds}
 \beta_\alpha=\frac{c_\alpha D_\alpha}{1+a_\alpha D_\alpha},\qquad
 \beta=\diag(\beta_1,\ldots,\beta_r),\qquad
 s_1=\Tr\beta,\quad s_2=\Tr\beta^2.
\end{equation}
Inactive coordinates have \(\beta_\alpha=0\). For all large \(r\), the
denominators are bounded away from zero, and
\begin{equation}\label{eq:insertion-coefficients}
 \|\beta\|\le C\theta,\qquad
 cT\le s_2\le CT,\qquad |s_1|\le CT
\end{equation}
with probability tending to one. The first and upper second bounds are
deterministic. The lower second bound follows from the positive proportion
of active roots, since \(|\beta_\alpha|\asymp\ell_r\) there.
For the last bound, pairing the two halves of each interval gives
\[
 |\EE\beta_\alpha|\le C\ell_r^2,
 \qquad \Var(\beta_\alpha)\le C\ell_r^2.
\]
Indeed the contributions \(\ell_r c(v)\) and
\(-\ell_r c(v+\ell_r)\) cancel to order \(\ell_r^2\), and the
factor \((1+aD)^{-1}\) differs from one by \(O(\ell_r)\).
Independence of the root environments now gives
\(|\EE s_1|\le CT\) and \(\Var s_1\le CT\).
Chebyshev's inequality makes the failure probability \(O(T^{-1})\);
the active-site count has exponentially small failure probability.
Along \(r=2^j\) these failures are summable, so
\eqref{eq:insertion-coefficients} holds eventually almost surely.

Let \(\mathcal E\) be the sigma field generated by \(\nu\) and all
sampled tree environments, including their root potentials. We shall
condition on almost every realization of \(\mathcal E\) for which
\eqref{eq:insertion-coefficients}, Lemma~\ref{lem:discrepancy}, and
Proposition~\ref{prop:transform-regularity} hold eventually.
After this conditioning, all coefficients and positions are deterministic;
the standardized Gaussian marks remain independent standard Gaussian variables.

Put \(I=\{j:|x_j|\le K\}\), \(m_I=|I|\), and
\(D_I=\diag(x_j:j\in I)\). In this section a \emph{raw GOE} matrix
has independent upper-triangular entries of variance one off the diagonal
and two on the diagonal, and centered Gaussian entries throughout.

\begin{proposition}[Gaussian insertion]\label{prop:gaussian-insertion}
On the preceding extension of any microscopic subsequential law, let
\(\mathcal W_r\) be the upper-left \(m_I\)-dimensional block of an
infinite raw GOE array independent of \(\mathcal E\), and set
\begin{equation}\label{eq:insertion-matrix}
 H_r=D_I+s_1\Id_{m_I}+\sqrt{s_2}\,\mathcal W_r.
\end{equation}
Then, along \(r=2^j\), the raw eigenvalue counting measures
\[
 \nu^{(r)}=\sum_{j=1}^{m_I}\delta_{\lambda_j(H_r)}
 \quad\Longrightarrow\quad\nu
\]
converge in law for the local vague topology, without conditioning.
Moreover \(m_I=2\rho K+O(K^{1-\chi})\) almost surely, and the
coefficient bounds \eqref{eq:insertion-coefficients} hold eventually
almost surely.
\end{proposition}

The proof follows the two cutoffs just described. First, truncation at
\(L\) turns the potential-switch determinant into the resolvent trace of
\(D_x+\mathcal G\beta\mathcal G^{\mathsf T}\), where \(D_x\) lists the
retained positions and \(\mathcal G\) has independent standard Gaussian
entries. Weighted resolvent estimates then remove the positions outside
\([-K,K]\). Finally, a second-order replacement of the independent
rank-one summands gives the scalar shift \(s_1\) and GOE variance \(s_2\)
in \eqref{eq:insertion-matrix}.

\subsection{The first cutoff and the finite matrix}

The discrepancy estimate and partial summation give
\begin{equation}\label{eq:sw-tail-sums}
 \nu([-R,R])=O(R),\qquad
 \left|\PV\sum_{|x_j|>R}\frac1{x_j}\right|=O(R^{-\chi}),\qquad
 \sum_{|x_j|>R}\frac1{x_j^2}=O(R^{-1}).
\end{equation}
These bounds hold for all sufficiently large \(R\), with constants
depending on \(\nu\). Write
\[
 M_r(z)=A_r+C_r^{1/2}
       \left(\PV\sum_j\frac{g_jg_j^\top}{x_j-z}\right)C_r^{1/2},
 \quad A_r=\diag(a_\alpha),\quad C_r=\diag(c_\alpha),
\]
using Proposition~\ref{prop:marked-resolvent}; the \(g_j\)'s are independent
standard Gaussian vectors of dimension \(r\). We restrict inverse
matrices below to active coordinates, whose number is at most \(r\).
All norm estimates through the first-cut comparison concern these active
compressions. On them \(C_r\) is bounded above and below uniformly;
inactive coordinates contribute the identity to the switch determinant.

For a fixed finite list of \(z\)'s in \(\HH\), let \(M_{r,L}\) be the
same expression restricted to \(|x_j|\le L\). Conditional on \(\mathcal E\),
\begin{equation}\label{eq:insertion-first-tail}
 \|M_r-M_{r,L}\|+\|M_r'-M_{r,L}'\|
       =O_{\PP}(r^2L^{-\chi}),\qquad
 \|M_r'\|=O_{\PP}(r^2).
\end{equation}
For the means use \eqref{eq:sw-tail-sums}. Each centered entry of the
tail has variance \(O(L^{-1})\), and summing squared entries proves
the bound in Hilbert--Schmidt norm; the derivative has an even smaller
tail. The derivative bound for the full matrix follows from
\(\sum_j|x_j-z|^{-2}<\infty\), first moments of Gaussian squares,
and the Hilbert--Schmidt or trace bound for its absolute summands.

With probability tending to one, both \(\Im M_r(z)\) and
\(\Im M_{r,L}(z)\) are at least \(r^{-5}\Id\) on active coordinates.
To see this, retain only points \(|x_j|\le r^4\). There are
\(2\rho r^4+O(r^{4(1-\chi)})\) such points; each Poisson weight is
at least \(c_zr^{-8}\). Their empirical Gaussian covariance tends
in operator norm to the identity: the expected squared Hilbert--Schmidt
error is \(O(r^2/r^4)=O(r^{-2})\). The uniform lower bound on \(C_r\)
then gives a lower bound of order \(r^{-4}\), and hence the claimed
weaker bound. The same retained points occur in the cutoff because
\(L>r^4\).

The potential increment matrix \(D=\diag(D_\alpha)\) is real and
invertible on active coordinates. Thus
\[
 \|(D^{-1}+M_r)^{-1}\|+
 \|(D^{-1}+M_{r,L})^{-1}\|\le2r^5
\]
on the preceding event. The inverse identity and
\eqref{eq:insertion-first-tail} therefore show that replacing \(M_r\)
by \(M_{r,L}\) in
\(\Tr((D^{-1}+M_r)^{-1}M_r')\) costs \(o_{\PP}(1)\).
Explicitly, after the trace and derivative are bounded, the inverse
difference multiplies the tail error by at most \(O_{\PP}(r^{13})\),
and \(r^{15}L^{-\chi}=r^{-85}\to0\).
The scalar principal-value transform \(m\) may also be cut at \(L\)
with an error tending to zero by \eqref{eq:sw-tail-sums}.

Let \(D_x=\diag(x_j:|x_j|\le L)\), and let \(\mathcal G\) have rows
indexed by these points and \(r\) independent standard Gaussian columns.
Determinant interchange gives
\begin{align*}
 \det(\Id+DM_{r,L}(z))
 &=\det(\Id+DA_r)\,
   \det\bigl(\Id+\beta\mathcal G^\top(D_x-z)^{-1}\mathcal G\bigr)\\
 &=\det(\Id+DA_r)\,
   \frac{\det(D_x+\mathcal G\beta\mathcal G^\top-z)}{\det(D_x-z)}.
\end{align*}
In particular, if
\[
 \mathcal S_r(z)=m(z)-\partial_z\log\det(\Id+DM_r(z)),
 \qquad A'=\mathcal G\beta\mathcal G^\top,
\]
then
\begin{equation}\label{eq:sw-finite-matrix}
 \mathcal S_r(z)=\Tr(D_x+A'-z)^{-1}+o_{\PP}(1).
\end{equation}
Only the logarithmic derivative is used, so no choice of a logarithm
branch is required. For every fixed \(r\), the joint unconditional law
of the imaginary parts of \(\mathcal S_r\) is that of \(m\), by
Proposition~\ref{prop:potential-switch}.

\subsection{Removing the outer positions}

The finite matrix in \eqref{eq:sw-finite-matrix} still has order \(L\).
We next compare its resolvent trace with that of the inner block indexed
by \(I\). This will leave a matrix of order \(K\), small enough for the
Gaussian replacement.

Let \(F=\{j:K<|x_j|\le L\}\), so that the indices of \(D_x\) split into \(I,F\).
For a fixed \(z\in\HH\), define the separate block resolvents
\[
 G_F=(D_F+A'_{FF}-z)^{-1},\qquad
 G_I=(D_I+A'_{II}-z)^{-1}.
\]
The upper-left block of the full resolvent is
\((G_I^{-1}-\Sigma)^{-1}\), where
\[
 \Sigma=A'_{IF}G_FA'_{FI}.
\]
Thus we must control both the direct far-block trace and the effect of
this Schur correction on the inner resolvent.

The diagonal positions give a natural weight on the far block, where
\(|x_j|>K\). On the inner block we regularize the same weight at the
larger height \(h_1=T^2\), where the perturbation tends to zero in the
weighted norm used below. Set
\begin{equation}\label{eq:sw-weights}
 h_1=T^2,\qquad U_I=(|D_I|+h_1)^{-1/2},\qquad
 U_F=|D_F|^{-1/2}.
\end{equation}
Partial summation using \eqref{eq:sw-tail-sums} gives
\begin{equation}\label{eq:sw-weight-bounds}
 \Tr U_I^2+\Tr U_F^2=O(\log r),\qquad
 \|U_I^2\|\le h_1^{-1},\quad \|U_F^2\|\le K^{-1}.
\end{equation}
For example, the inner sum is bounded by integrating the linear count
bound against \((|x|+h_1)^{-1}\); the far sum is treated in the same way
between \(K\) and \(L\). Both endpoints are fixed powers of \(r\).
In the following estimates, \(\operatorname{polylog}(r)\) denotes a fixed
power of \(\log r\), whose exponent may increase from one occurrence
to the next.

We will show that \(G_F\) is close to \(D_F^{-1}\) in these weights,
that \(\|U_I\Sigma U_I\|=o_{\PP}(T^{-5})\), and that
\(\|U_I^{-1}G_I(z)U_I^{-1}\|=O_{\PP}(T^2)\).
The last two estimates make the inner resolvent change small even
after the two inverse factors in the resolvent identity. The logarithmic
trace bounds above then turn weighted norm control into trace control.
We begin with the Gaussian block estimates that supply these bounds.

\begin{lemma}[Weighted block estimates]\label{sw:weighted-blocks}
For \(\mathcal A=I,F\), put \(a_I=h_1\), \(a_F=K\). With probability tending
to one,
\begin{equation}\label{eq:sw-bernstein}
 \left\|U_{\mathcal A}
       (A'_{\mathcal A\mathcal A}-s_1\Id)
       U_{\mathcal A}\right\|
 \le \operatorname{polylog}(r)
       \left(\sqrt{T/a_{\mathcal A}}+\theta\right).
\end{equation}
Moreover, for \(P_0=U_IA'_{IF}U_F\),
\begin{equation}\label{eq:sw-cross-block}
 \|P_0\|_{\HS}^2\le T\operatorname{polylog}(r)
\end{equation}
with probability tending to one.
\end{lemma}

\begin{proof}
The centered matrix in \eqref{eq:sw-bernstein} is a sum of independent
self-adjoint matrices
\[
 \sum_{\alpha=1}^{r}\beta_{\alpha\alpha}
       (v_\alpha v_\alpha^\top-P),\qquad P=U_{\mathcal A}^2,
\]
where \(v_\alpha\) is Gaussian of covariance \(P\).
Its coefficients may have either sign. Gaussian fourth moments give
\begin{equation}\label{eq:sw-fourth-moment}
 \EE[(v_\alpha v_\alpha^\top)^2]
   =(\Tr P)P+2P^2.
\end{equation}
Truncate each original Gaussian column on the event that all its entries
have absolute value at most \(\log r\), and center the truncated summand.
There are only polynomially many entries, so this alters the matrices with
superpolynomially small probability. Gaussian tails also make the resulting
centering correction superpolynomially small.

Each truncated summand has norm at most
\[
 C\|\beta\|(\log r)^2\Tr P
 \le \theta\operatorname{polylog}(r).
\]
If \(X\) is a self-adjoint random matrix, then
\(\EE[(X-\EE X)^2]=\EE[X^2]-(\EE X)^2\le\EE[X^2]\).
Also truncating \(v_\alpha v_\alpha^\top\) by an event can only decrease
the expectation of its square in positive semidefinite order.
Thus \eqref{eq:sw-fourth-moment}, including for negative
\(\beta_{\alpha\alpha}\), bounds the matrix variance parameter by
\[
 C s_2(\Tr P)\|P\|
 \le \frac{T}{a_{\mathcal A}}\operatorname{polylog}(r).
\]
The bounded self-adjoint matrix Bernstein inequality
\cite[Theorem~1.4]{Tropp2012}, applied to both signs, gives, for a sum of independent centered
self-adjoint matrices of norm at most \(b\) and variance parameter
\(\sigma_{\rm mat}^2\), a two-sided tail bounded by
\[
 2|\mathcal A|\exp\left(
 -\frac{c t^2}{\sigma_{\rm mat}^2+bt}\right).
\]
Its dimension factor is polynomial in \(r\). Taking a sufficiently large
power of \(\log r\) in \(t\) proves \eqref{eq:sw-bernstein}.

For the cross block, independence of distinct Gaussian rows and direct
second-moment calculation give
\[
 \EE\|P_0\|_{\HS}^2
 =s_2(\Tr U_I^2)(\Tr U_F^2).
\]
Equation~\eqref{eq:sw-cross-block} follows from
\eqref{eq:sw-weight-bounds} and Markov's inequality, with one additional
logarithmic factor.
\end{proof}

We first apply the weighted block estimates to the far resolvent
\(G_F\) and the Schur correction \(\Sigma\), keeping \(z\in\HH\) fixed.

\begin{lemma}[Far block and Schur correction]\label{sw:schur}
Let
\[
 \delta_F=\operatorname{polylog}(r)
       \left(\frac{1+T}{K}+\sqrt{\frac TK}+\theta\right).
\]
Then
\begin{equation}\label{eq:sw-far-resolvent}
 \|U_F^{-1}(G_F-D_F^{-1})U_F^{-1}\|\le\delta_F=o(1)
\end{equation}
with probability tending to one, and \(\Tr G_F=o_{\PP}(1)\).
For the Schur correction \(\Sigma\),
\begin{equation}\label{eq:sw-schur-small}
 \|U_I\Sigma U_I\|=o_{\PP}(T^{-5}).
\end{equation}
\end{lemma}

\begin{proof}
Conjugating the matrix to be inverted by \(U_F\) gives
\[
 U_F(D_F+A'_{FF}-z)U_F
 =\operatorname{sign}(D_F)+U_F(A'_{FF}-z)U_F.
\]
The second term has norm at most \(\delta_F\), by
Lemma~\ref{sw:weighted-blocks}, the bound \(|s_1|=O(T)\), and
\eqref{eq:sw-weight-bounds}. Since \(\delta_F=o(1)\), a Neumann expansion
proves \eqref{eq:sw-far-resolvent}. In particular,
\(Q_F=U_F^{-1}G_FU_F^{-1}\) has bounded norm.
By \eqref{eq:sw-tail-sums},
\[
 \left|\sum_{j\in F}\frac1{x_j}\right|=O(K^{-\chi}),
 \qquad
 |\Tr(G_F-D_F^{-1})|
 \le\delta_F\Tr U_F^2=o(1).
\]
This proves the assertion about the far trace.

For the Schur correction, replacing \(G_F\) by \(D_F^{-1}\) costs at most
\begin{equation}\label{eq:sw-schur-replacement}
 \|P_0\|_{\HS}^2\delta_F
 \le T\operatorname{polylog}(r)\delta_F
\end{equation}
in the norm weighted on both sides by \(U_I\).
It remains to estimate
\(A'_{IF}D_F^{-1}A'_{FI}\).
Condition further on the inner Gaussian row block \(\mathcal G_I\), and set
\[
 Q=\mathcal G_I\beta^2\mathcal G_I^\top.
\]
The columns of \(A'_{IF}\), indexed by \(j\in F\), are then independent
centered Gaussian vectors \(a_j\) with covariance \(Q\).
Consequently
\begin{align*}
 \EE\left[\sum_{j\in F}\frac{a_ja_j^\top}{x_j}
       \,\middle|\,\mathcal G_I,\mathcal E\right]
 &=Q\sum_{j\in F}\frac1{x_j},\\
 \Var\left(\sum_{j\in F}\frac{(a_j)_a(a_j)_b}{x_j}
       \,\middle|\,\mathcal G_I,\mathcal E\right)
 &=(Q_{aa}Q_{bb}+Q_{ab}^2)
       \sum_{j\in F}\frac1{x_j^2}.
\end{align*}
Also
\(\EE[Q_{aa}^2\mid\mathcal E]
 =s_2^2+2\Tr\beta^4=O(T^2)\).
Since \(Q\) is positive semidefinite,
\(|Q_{ab}|^2\le Q_{aa}Q_{bb}\); Cauchy--Schwarz bounds the remaining
second moments by \(O(T^2)\).
Multiplying entry \(a,b\) by \((U_I)_{aa}(U_I)_{bb}\), summing its second
moment over \(a,b\), and using \eqref{eq:sw-tail-sums} and
\eqref{eq:sw-weight-bounds}, we obtain
\[
 \EE\left\|
 U_IA'_{IF}D_F^{-1}A'_{FI}U_I
 \right\|_{\HS}^2
 \le C T^2(\log r)^2(K^{-2\chi}+K^{-1}).
\]
Markov's inequality therefore bounds this norm by
\(T\operatorname{polylog}(r)(K^{-\chi}+K^{-1/2})\) with probability
tending to one.

Combining this estimate with \eqref{eq:sw-schur-replacement} gives
\[
 \|U_I\Sigma U_I\|
 \le \operatorname{polylog}(r)\,O\left(
 TK^{-\chi}+TK^{-1/2}+\frac{T(1+T)}K
       +T\sqrt{T/K}+T\theta\right)
\]
with probability tending to one. After multiplication by \(T^5\), all
terms tend to zero by a fixed power of \(r\). For example, the three
potentially slow terms have exponents
\[
 6\tau-\frac{\chi}{100},\qquad
 \frac{13}{2}\tau-\frac1{200},\qquad
 \frac{13}{2}\tau-\frac12,
\]
respectively. They are negative under \eqref{eq:sw-parameters}; the terms
\(T^6K^{-1/2}\) and \(T^6(1+T)/K\) have still ample margins.
This proves \eqref{eq:sw-schur-small}.
\end{proof}

We now convert the small weighted Schur correction into a small trace
error. The fixed-height inner resolvent cannot be estimated directly from
\eqref{eq:sw-bernstein}; instead we first work at the larger height \(h_1=T^2\).

\begin{lemma}[Inner resolvent and trace cut]\label{sw:trace-cut}
For every fixed \(z\in\HH\),
\begin{equation}\label{eq:sw-inner-resolvent}
 \|U_I^{-1}G_I(z)U_I^{-1}\|=O_{\PP}(h_1),
\end{equation}
and
\begin{equation}\label{eq:sw-second-cut}
 \Tr(D_x+A'-z)^{-1}
 =\Tr(D_I+A'_{II}-z)^{-1}+o_{\PP}(1).
\end{equation}
The same estimates hold jointly for a fixed finite list of \(z\)'s.
\end{lemma}

\begin{proof}
Every diagonal entry of \(U_I(D_I-ih_1)U_I\) has modulus at least
\(1/\sqrt2\). The added term has norm at most
\[
 \|U_IA'_{II}U_I\|
 \le C T/h_1+\operatorname{polylog}(r)
       \left(\sqrt{T/h_1}+\theta\right)=o(1)
\]
with probability tending to one. Hence
\[
 \|U_I^{-1}G_I(ih_1)U_I^{-1}\|\le C.
\]
The identity \(G_I(ih_1)^*G_I(ih_1)=h_1^{-1}\Im G_I(ih_1)\) now gives
\[
 \|G_I(ih_1)U_I^{-1}\|+\|U_I^{-1}G_I(ih_1)\|
 \le C h_1^{-1/2}.
\]
Alternatively the same mixed-weight estimates follow from the preceding
weighted inverse and \(\|U_I\|\le h_1^{-1/2}\).
Use the exact twice-iterated resolvent identity
\begin{align*}
 G_I(z)
 ={}&G_I(ih_1)+(z-ih_1)G_I(ih_1)^2\\
 &+(z-ih_1)^2G_I(ih_1)G_I(z)G_I(ih_1).
\end{align*}
The unweighted middle resolvent in the last term has norm at most
\((\Im z)^{-1}\). Conjugating the identity by \(U_I^{-1}\), the three
terms are bounded by
\[
 C,\qquad C|z-ih_1|/h_1,\qquad
 C|z-ih_1|^2/(h_1\Im z),
\]
respectively. This proves \eqref{eq:sw-inner-resolvent}.

The Schur complement formula for the full inverse has upper-left block
\(\widetilde G_I=(G_I^{-1}-\Sigma)^{-1}\).
Set \(R_I=U_I^{-1}G_IU_I^{-1}\) and \(S_I=U_I\Sigma U_I\).
Then
\[
 U_I^{-1}\widetilde G_IU_I^{-1}
 =(R_I^{-1}-S_I)^{-1}.
\]
By \eqref{eq:sw-inner-resolvent} and
\eqref{eq:sw-schur-small}, a Neumann expansion gives
\[
 \|U_I^{-1}\widetilde G_IU_I^{-1}\|=O_{\PP}(h_1),\qquad
 \|U_I^{-1}(\widetilde G_I-G_I)U_I^{-1}\|
 \le O_{\PP}(h_1^2)o_{\PP}(T^{-5})=o_{\PP}(T^{-1}).
\]
Multiplying the last bound by \(\Tr U_I^2=O(\log r)\) shows that their
traces differ by \(o_{\PP}(1)\).

The full block-inverse trace is
\[
 \Tr\widetilde G_I+\Tr G_F+
 \Tr\bigl(\widetilde G_IA'_{IF}G_F^2A'_{FI}\bigr).
\]
The far trace is \(o_{\PP}(1)\) by Lemma~\ref{sw:schur}.
For the last term recall \(Q_F=U_F^{-1}G_FU_F^{-1}\).
It has bounded norm, and
\[
 U_F^{-1}G_F^2U_F^{-1}=Q_FU_F^2Q_F,\qquad
 \|Q_FU_F^2Q_F\|=O_{\PP}(K^{-1}).
\]
No positivity is asserted for this complex square. The ordinary
operator-norm and Hilbert--Schmidt trace inequality instead gives
\begin{align*}
 \left|\Tr\bigl(\widetilde G_IA'_{IF}G_F^2A'_{FI}\bigr)\right|
 &\le
 \|U_I^{-1}\widetilde G_IU_I^{-1}\|\,
 \|P_0\|_{\HS}^2\,
 \|Q_FU_F^2Q_F\|\\
 &\le \operatorname{polylog}(r)\,O_{\PP}(T^3/K)=o_{\PP}(1).
\end{align*}
This proves \eqref{eq:sw-second-cut}. All estimates used the same Gaussian
matrices, so taking a finite union gives the joint assertion.
\end{proof}

\subsection{Gaussian replacement and convergence of point measures}

We have reduced the switched trace to an inner matrix of size
\[
 m_I=2\rho K+O(K^{1-\chi})\asymp K.
\]
Its perturbation is a sum of \(r\) independent small rank-one matrices.
Only at this reduced size do we replace that sum by GOE noise.

\begin{lemma}[Gaussian replacement]\label{sw:lindeberg}
Conditional on almost every realization of \(\mathcal E\), bounded smooth tests with bounded
derivatives of finitely many trace values have expectation difference
\(o(1)\) between
\[
 D_I+A'_{II}
 \quad\hbox{and}\quad
 H_r=D_I+s_1\Id_{m_I}+\sqrt{s_2}\,\mathcal W.
\]
\end{lemma}

\begin{proof}
Write
\[
 A'_{II}-s_1\Id_{m_I}
 =\sum_{\alpha=1}^{r}
 \beta_{\alpha\alpha}(g_\alpha g_\alpha^\top-\Id_{m_I}),
\]
where the \(g_\alpha\) are independent standard \(m_I\)-dimensional
Gaussian vectors. The covariance identity
\[
 \Cov(g_ag_b,g_cg_d)
 =\delta_{ac}\delta_{bd}+\delta_{ad}\delta_{bc}
\]
shows that \(gg^\top-\Id\) and a raw GOE matrix have the same entry
covariances. Their means are both zero. Thus replace the summands one at
a time by \(\beta_{\alpha\alpha}\mathcal W_\alpha\), with independent raw
GOE matrices \(\mathcal W_\alpha\). Negative coefficients are allowed:
both covariances are multiplied by \(\beta_{\alpha\alpha}^2\).

For either kind of summand \(X\),
\[
 \EE\|X\|^3\le C\theta^3m_I^3.
\]
For the rank-one term this follows from Gaussian moments of \(\|g\|^2\);
for the GOE term the Hilbert--Schmidt norm gives a sufficient bound.
Fix spectral parameters \(z_1,\ldots,z_b\in\HH\), and a bounded smooth
test \(\Phi\) on the corresponding real and imaginary trace coordinates,
with bounded derivatives. For any real symmetric matrix \(A\), its
resolvents satisfy \(\|(A-z_\ell)^{-1}\|\le(\Im z_\ell)^{-1}\).
Differentiating the resolvent in a real symmetric direction \(\Delta\)
gives, for \(j=1,2,3\), a trace derivative bounded by
\(C_{z,j}m_I\|\Delta\|^j\). The chain rule therefore bounds the third
directional derivative of the composed test by
\[
 C_{z,\Phi}m_I^3\|\Delta\|^3.
\]
This estimate holds at every intermediate real symmetric matrix,
independently of eigenvalue gaps.

The Lindeberg replacement argument (see \cite{Chatterjee2006}) uses Taylor expansion to second order, which cancels the means and covariances in
each replacement. Summing the third-order remainders gives total error
\[
 C_{z,\Phi}m_I^6 r\theta^3
 =O\bigl(r^{-0.44+(3/2)\tau}\bigr)=o(1).
\]
Finally \(\sum_\alpha\beta_{\alpha\alpha}\mathcal W_\alpha\) has exactly
the law \(\sqrt{s_2}\,\mathcal W\), proving the lemma.
\end{proof}

\begin{proof}[Completion of Proposition~\ref{prop:gaussian-insertion}]
The coefficient bounds were proved in \eqref{eq:insertion-coefficients}.
Combine \eqref{eq:sw-finite-matrix} with the trace cut in
Lemma~\ref{sw:trace-cut} and the Gaussian replacement in
Lemma~\ref{sw:lindeberg}. They give a conditional comparison of bounded
smooth tests of finitely many trace values with \(\mathcal S_r\).
Its error tends to zero for almost every \(\mathcal E\), and is bounded,
so bounded convergence removes the conditioning. The unconditional
potential-switch identity from Proposition~\ref{prop:potential-switch}
therefore gives
\begin{equation}\label{eq:sw-trace-limit}
 \bigl(\Im\Tr(H_r-z_\ell)^{-1}\bigr)_{\ell=1}^b
   \ \Longrightarrow\
 \bigl(\Im m(z_\ell)\bigr)_{\ell=1}^b
 \qquad(z_1,\ldots,z_b\in\HH).
\end{equation}
Bounded smooth tests determine finite-dimensional weak convergence.
In particular, switching invariance has only been used after averaging
out the conditioning data.

We finish by explaining why these trace limits imply convergence of point
measures, without assuming a uniform bound on their total number of
points. Let \(\nu^{(r)}\) be the eigenvalue counting measure of \(H_r\), and
define
\[
 \sigma_r(\,d x)=\frac{\nu^{(r)}(\,d x)}{1+x^2}.
\]
Its total mass is \(\Im\Tr(H_r-i)^{-1}\), whose laws are tight by
\eqref{eq:sw-trace-limit}. Regarding \(\sigma_r\) as a finite positive
measure on the one-point compactification
\(\overline{\RR}=\RR\cup\{\infty\}\), its laws are therefore tight.
For \(z=e+i\eta\in\HH\), the kernel
\[
 k_z(x)=(1+x^2)\frac{\eta}{(x-e)^2+\eta^2}
\]
extends continuously to \(\overline{\RR}\), with \(k_z(\infty)=\eta\).
Its integral against \(\sigma_r\) is \(\Im\Tr(H_r-z)^{-1}\).
It follows that in any subsequential limit \(\sigma\), the joint law of
these integrals at a countable dense set of \(z\)'s is that of
\(\Im m(z)\). Continuity extends this conclusion to all \(z\).

These integrals determine the de-weighted restriction of \(\sigma\) to
\(\RR\). Indeed, for \(f\in C_c(\RR)\), integrate the kernel against
\(\pi^{-1}f(e)\,d e\) and let \(\eta\downarrow0\). On \(\RR\) the resulting
kernel tends to \((1+x^2)f(x)\), by the Poisson approximate identity.
The integrated kernels are bounded uniformly over
\(x\in\overline{\RR}\) and \(0<\eta\le1\): near the compact support this is
the Poisson mass bound, and away from it the factor \(1+x^2\) cancels
the quadratic denominator. At infinity their value is
\(\eta\pi^{-1}\int f(e)\,d e\), which tends to zero.
Dominated convergence is thus applicable to each finite measure
\(\sigma\), including any mass at infinity.

By Proposition~\ref{prop:marked-resolvent}, applying the same inversion to
\(\Im m(z)\) gives \(\int f\,\,d\nu\). Taking a countable determining
family of compactly supported tests, the joint distribution of these
integrals is therefore exactly the law of \(\nu\).
The map from \(\sigma\) to its de-weighted restriction is continuous for
local vague convergence: for a compactly supported \(f\), the function
\((1+x^2)f(x)\), extended by zero at infinity, is continuous on
\(\overline{\RR}\). Every subsequential local vague limit of
\(\nu^{(r)}\) consequently has the law of \(\nu\).
This proves the point-process assertion.
\end{proof}
